% New prospective section. Not assembled, rendered, published or scientifically accepted.
% Preserve earlier studies, result tables, citations and frozen failures as historical evidence.
\section{Measure both ends. Test the translation.}

The following prospective protocol tests the paired learner described in ``Two pairs. One conditional learner'': can permitted observations on one side of a sensorimotor task predict withheld measurements on the other, on recordings the model has not seen? The apparatus remains non-invasive EEG, non-invasive surface EMG, a calibrated motion-capture relay-composite, an open musculoskeletal-simulation physics-engine with a lifelike explanatory view, and an open multimodal model. Its acquisition and validation are unfinished.

The earlier brain--muscle and measured-movement--virtual-body comparisons define the two pairs. Here the protocol specifies their records, targets, calibration and independent tests. Anatomical form constrains possible mechanisms; measured function must test the proposed connections. A conditional learner's combined prediction must answer to measurements outside itself.

\subsection{Keep the observations distinct}

Keep raw EEG, raw EMG, derived features and movement estimates in separate
records. Identify the participant and session in each record, then declare the
sensor sites, references, coordinate conventions, units, sampling and timestamp
semantics, and clocks. Trace every processing step and preserve missing intervals and uncertainty. Retain task
instructions, stimulus events, observed behaviour and any separately obtained
participant report. Independent participants or recordings cannot supply a
simultaneous brain--muscle--movement observation.

Surface EMG records muscle electrical activity. A sensory-loop claim also needs
known stimuli and appropriate measured feedback, contact or force channels;
muscle recordings do not provide incoming sensory-nerve traffic. Preserve the
actual spatial resolution of EEG measurements and source estimates. Templates
and explanatory cutaways do not register participants or localize individual
neurons.

\subsection{The spine shapes the loop}

The spinal cord changes how sensory input becomes muscle output. Account for
local coordination, sensory gating, reflex gain and some lasting adaptations,
rather than reducing it to a transmission delay. Separate a response to the
current posture or task from a change retained after training. Keep both apart
from adaptation in sensory receptors, peripheral pathways, neuromuscular
junctions and mechanical tissues. These processes interact; changed output alone
does not locate a memory.

Human H-reflex conditioning and the constrained mouse adaptation studied by
Lavaud and colleagues ask different questions.\cite{thompson2009retained,thompson2013retained,lavaud2024}
The human experiments do not establish brain-independent skill learning. The
mouse study does not establish unrestricted walking or reaching. Circuits
required for acquisition or recall are not thereby exclusive memory-storage
sites. Our human comparison uses the abstracts, selected results and limited
methods passages. The mouse comparison also draws on relevant author-manuscript
passages; equivalence to the final published edition and the complete methods
and statistics have not been established here. Passive proprioceptive training
motivates a peripheral observation model; its psychophysical and spindle-afferent
experiments were separate.\cite{ackerley2022retained} Those observations do not
establish complex memory stored in peripheral axons.

Define the function before choosing its computational substitute. Candidate
functions include gain and inhibition, phase- and load-dependent sensory
responses, coordination across supported muscles or joints, relevant rhythms,
and acquisition, retention and expression of a declared adaptation. Give each
function its own task, species, anatomical scope, measurement and reference.
Naming a mechanism does not show that the model implements it. A hand task
cannot inherit an unmeasured locomotor circuit.

Record task context and history. Distinguish the state within a trial from a
learned parameter change; retain elapsed time, intervening activity, resets,
drift and uncertainty. Recorded spinal or peripheral activity needs supported
channels of its own. Where those observations are absent, label internal state
as estimated or as a simulated controller. EEG, EMG and movement cannot be
relabeled as recorded spinal activity or identified interneurons.

Declare controls and endpoints before fitting or experimentation. Test
immediate context effects separately from retention after a declared interval
and from peripheral or tissue changes. Preserve the original held-out subjects
and tasks, eligible independent measurements, matched baselines and ablations.
Check alternative explanations, time shifts, missing inputs and artifact
leakage. Lower movement error alone cannot locate a biological adaptation,
establish its cause or certify useful plasticity.

To externalize a specified function, propose a substitute that can be tested
through declared peripheral and cerebral interfaces. Compare latency, state
dependence, stability, uncertainty and measured outcomes with the specified
biological function. Computational feedback, peripheral stimulation and restored
biological function require different evidence. A matched virtual trajectory
does not demonstrate spinal replacement, relocated learning or clinical repair.
Keep the spinal-cord site, measured function and proposed computational
substitute distinct. Physical spinal scope remains read-only.

\subsection{Give the two directions different targets}

For brain-to-limb prediction, give the model EEG and only the declared context
available at inference. Test its predicted muscle and movement observables
against withheld measurements. A brain-only condition receives no measured EMG
or movement input. If those inputs help, evaluate them in a separately labeled
hybrid condition.

For limb-to-brain prediction, give the model muscle activity, movement and
declared context. Its answer must allow for different patterns in the supported
recorded EEG features. A conditional distribution assigns probabilities to
specified feature values or ranges, given those inputs. Compare that prediction
with independently retained EEG, then report uncertainty and alternatives.
Movement can narrow the possibilities without identifying a unique hidden
brain state or subjective experience. Agreement between learned models is a
consistency check; it cannot supply the measured target.

\subsection{Reading a signal and writing an output}

The protocol records EEG and EMG; its two prediction directions do not themselves deliver sensation or enact a movement. ``Write sensation. Enact action'' defines the computer and physiological boundaries. Testing a writer would additionally require a separately qualified sensory writer or transducer, a specified actual response and distinct device, measurement, safety and clinical evidence. Indirect sensory feedback and direct neural writing remain different interfaces; a predicted EEG feature is neither. Enaction needs an appropriate virtual, robotic or biological output interface. A robot command has an action role, not the anatomical role of an efferent nerve.

Physical sensory writers, transducers and body actuators remain excluded from the acquisition and virtual-work plan, unqualified and unpriced. This proposal authorizes no stimulation, participant acquisition, biological actuation, new equipment purchase or training.

\subsection{Align before learning}

Calibrate the acquisition views. Verify their clocks against recorded events,
offsets and drift, while preserving occlusion and dropout. Keep a missing
observation missing: a predicted finger position is not a measurement. Declare
frames, scale, joint conventions and acceptable spatial, temporal and end-to-end
errors before selecting equipment or models. Position, joint, contact and force
claims each need an appropriate independent reference.

Specify how the virtual body represents subject scale, muscles, tendons, joint
constraints and contact. Declare its integration settings and observation
models. Biomechanical simulation predicts mechanics; VFX rendering makes those
predictions visible. Judge individual parts and their assembly against the
original anatomical, physical, functional and real-reference-series criteria.

\subsection{Build an open corpus without inventing pairs}

Inventory neural structure, neural activity, muscle structure and muscle
activity as separate domains. Use measured movement and task context as
constraints. For each source, record rights, version, access conditions,
species, participants, tasks, montages, units, clocks, preprocessing and missing
paired modalities. Broad open-data coverage remains a continuing objective,
not an assertion that every dataset is accessible, licensed, compatible or
already used. Check code, pretrained weights and training-data rights separately.

The candidate uses an encoder for each modality, explicit input-availability
masks, a shared conditional representation and separate decoders for the two
directions. Keep simple EEG-only, EMG-only, movement-only and hybrid baselines.
Define pretraining, post-training and fine-tuning tasks before fitting. An
existing encoder supplies a candidate component, not an already established
brain--body model.

\subsection{Make the reward answer to measurements}

First test paired supervised prediction in each direction. Compare any later\linebreak[4]
reinforcement-learning or other closed-loop refinement with that baseline.
Evaluate independent movement measurements, task outcomes, relevant contact or
force, latency and physical validity. Freeze reward units, terms, weights and
decision thresholds before evaluation. The candidate cannot invent the target
against which it earns a reward.

Backpropagation updates model parameters. Reverse prediction estimates another
observable. Measured biological feedback comes from the body. Identify which
process the study actually implements. A supervised prediction error alone does
not establish reinforcement learning: a proposed control study must also declare
its state, action, transition and reward.

Partition eligible participants, sessions and novel tasks before preprocessing
or model selection. Preserve TRAIN-only normalization, frozen seeds and
checkpoint selection, untouched held-out evaluation, uncertainty and original
negative controls. Include informative missing-modality, time-shift, artifact,
drift and alternative-inverse comparisons. Keep reserved recordings closed and
retain completed fits and failed benefit conjunctions. Evaluate frozen inference
separately from any proposed adaptive update. This prospective protocol admits
no new fitting, optimizer update or model/backend promotion.

\subsection{Show what was measured and what was inferred}

Show the selected source clock and the inputs available to each prediction.
Place recorded EEG, muscle activity and movement beside their predictions, with
uncertainty, errors, gaps and stimulus or contact context. Identify recordings,
inferred features and simulation as separate layers. Preserve site/channel
inspection and source-time play, pause, seek and step. Unsupported peripheral
or spinal pathways remain visible gaps. An animated cord cannot substitute for
a spinal recording or supported anatomical localization.